\documentclass[a4paper,11pt]{article}
\usepackage[utf8]{inputenc}
\usepackage[french]{babel}
\usepackage[T1]{fontenc}
\usepackage{amsmath,amssymb,amsthm}
\usepackage{geometry}
\geometry{margin=2cm}
\usepackage{enumitem}
\usepackage{hyperref}
\usepackage{booktabs}
\usepackage{array}
\usepackage{xcolor}
\usepackage{listings}
\usepackage{titlesec}
\usepackage{fancyhdr}
\usepackage{lastpage}
\usepackage{graphicx}
\usepackage{etoolbox}
\usepackage{multirow}
\AtBeginDocument{\shorthandoff{;:!?}}

\titleformat{\section}{\large\bfseries\color{blue!60!black}}{}{0pt}{}[\vspace{-0.5ex}\rule{\textwidth}{0.4pt}]
\titleformat{\subsection}{\bfseries\color{blue!40!black}}{}{0pt}{}
\titleformat{\subsubsection}{\itshape\color{blue!30!black}}{}{0pt}{}

\lstdefinestyle{monstyle}{
  frame=single,
  numbers=left,
  numbersep=5pt,
  breaklines=true,
  basicstyle=\small\ttfamily,
  backgroundcolor=\color{gray!5},
  rulecolor=\color{gray!40},
  columns=fullflexible,
  keepspaces=true,
  showstringspaces=false
}
\lstset{style=monstyle}

\pagestyle{fancy}
\fancyhf{}
\fancyhead[L]{\small STA211 -- R\'eseaux de neurones (cours 2)}
\fancyhead[R]{\small Fiche r\'ecapitulative}
\fancyfoot[C]{\thepage/\pageref{LastPage}}
\renewcommand{\headrulewidth}{0.4pt}

\theoremstyle{definition}
\newtheorem*{definition}{D\'efinition}
\theoremstyle{remark}
\newtheorem*{remark}{Remarque}
\newtheorem*{important}{\`A retenir}

\title{\textbf{Fiche r\'ecapitulative -- STA211}\\[4pt]
\large R\'eseaux de neurones -- Cours 2 \\[2pt]
\small R\'eseaux de neurones convolutionnels (ConvNets) et LeNet}
\author{Nicolas Thome (CNAM)}
\date{14 juin 2026}

\begin{document}
\maketitle
\thispagestyle{empty}

\begin{abstract}
Cette fiche pr\'esente les r\'eseaux de neurones convolutionnels
(ConvNets / CNN), qui tirent parti de la structure spatiale des
donn\'ees (images, audio) via la convolution, le partage de poids
et le pooling. Le mod\`ele historique LeNet-5 (LeCun et al., 1989)
est d\'etaill\'e comme \'etude de cas avec son architecture
compl\`ete et ses performances sur MNIST.
\end{abstract}

\vfill
\tableofcontents
\newpage

% ===================================================================
\section{Limites des r\'eseaux fully connected (MLP)}
% ===================================================================

Les r\'eseaux enti\`erement connect\'es (MLP) pr\'esentent trois
limitations majeures pour les donn\'ees de signaux~:

\begin{enumerate}
  \item \textbf{Explosion du nombre de param\`etres}~: m\^eme une
    seule couche cach\'ee g\'en\`re un nombre de poids colossal
    (ex~: image $32\times 32$ en entr\'ee, couche cach\'ee de $1000$
    neurones $\Rightarrow$ $1$ million de param\`etres).
  \item \textbf{Ignorance de la topologie spatiale}~: la
    repr\'esentation vectorielle perd toute notion de voisinage.
    Le MLP est indiff\'erent aux permutations des pixels d'entr\'ee.
  \item \textbf{Absence d'invariance}~: un faible d\'ecalage en
    entr\'ee peut produire une repr\'esentation compl\`etement
    diff\'erente.
\end{enumerate}

% ===================================================================
\section{La convolution}
% ===================================================================

\subsection{Connectivit\'e locale et partage de poids}

Pour pallier ces limites, on introduit deux id\'ees cl\'es~:

\begin{enumerate}
  \item \textbf{Connectivit\'e locale (sparse connectivity)}~:
    chaque neurone cach\'e n'est connect\'e qu'\`a une petite r\'egion
    locale de l'entr\'ee (champ r\'ecepteur), inspir\'e du cortex
    visuel.
  \item \textbf{Partage de poids (shared weights)}~: le m\^eme filtre
    (noyau) est appliqu\'e \`a toutes les positions spatiales,
    r\'eduisant consid\'erablement le nombre de param\`etres.
\end{enumerate}

\subsection{Convolution 2D}

Soit une image $f$ et un filtre $h$ de taille impaire $d$. La
convolution 2D est d\'efinie par~:
\begin{equation}
  f'(i,j) = (f \star h)(i,j) =
  \sum_{n=-\frac{d-1}{2}}^{\frac{d-1}{2}}
  \sum_{m=-\frac{d-1}{2}}^{\frac{d-1}{2}}
  f(i-n, j-m)\,h(n,m)
\end{equation}

\noindent La sortie pour $1$ filtre est une carte 2D~; pour $K$
filtres on obtient $K$ cartes 2D.

\subsection{Convolution vs corr\'elation crois\'ee}

La corr\'elation crois\'ee (\emph{cross-correlation}) est
similaire mais sans sym\'etrisation du masque~:
\begin{equation}
  f'(i,j) = (f \otimes h)(i,j) =
  \sum_{n}\sum_{m} f(i+n, j+m)\,h(n,m)
\end{equation}
En pratique, les frameworks Deep Learning impl\'ementent la
corr\'elation crois\'ee (plus simple) et les filtres sont appris.

\subsection{Convolution sur tenseurs (couche convolutionnelle)}

Pour une image couleur (3 canaux) ou une carte de features
profonde, la convolution s'applique sur toute la profondeur~:
\begin{equation}
  f'(i,j) = \sum_{k=1}^{K} \sum_{n}\sum_{m}
  f(i-n, j-m, k)\,h(n,m,k) + b
\end{equation}

\noindent Chaque filtre produit une carte 2D~; $F$ filtres
produisent $F$ cartes empil\'ees en un \textbf{tenseur} de
profondeur $F$.

\begin{important}
  La couche convolutionnelle est une op\'eration affine
  (convolution lin\'eaire + biais), suivie d'une non-lin\'earit\'e
  (ReLU, sigmo\"ide, etc.).
\end{important}

% ===================================================================
\section{Le pooling}
% ===================================================================

\subsection{Objectif}

Le \textbf{pooling} (sous-\'echantillonnage) permet de~:
\begin{itemize}
  \item R\'eduire la r\'esolution spatiale.
  \item Apporter une invariance aux translations locales.
  \item Diminuer le nombre de param\`etres.
\end{itemize}

\subsection{Max pooling}

La valeur maximale dans une fen\^etre $p \times p$ est retenue avec
un pas (stride) $s$~:
\begin{itemize}
  \item Exemple~: fen\^etre $2\times 2$, stride $2$ $\Rightarrow$
    division par $2$ de chaque dimension spatiale.
  \item Offre une invariance partielle aux translations~: un
    d\'ecalage inf\'erieur \`a la taille de la fen\^etre de pooling
    ne modifie pas n\'ecessairement le maximum.
\end{itemize}

\subsection{Types de pooling}

\begin{tabular}{ll}
\toprule
Type & Op\'eration \\
\midrule
Max pooling & Maximum dans la fen\^etre \\
Average pooling & Moyenne dans la fen\^etre \\
Sum pooling & Somme (pond\'er\'ee par un param\`etre appris) \\
\bottomrule
\end{tabular}

% ===================================================================
\section{Architecture d'un ConvNet}
% ===================================================================

Un r\'eseau convolutionnel (ConvNet / CNN) est constitu\'e de
l'empilement de blocs \'el\'ementaires~:
\begin{equation}
  \boxed{\text{Convolution} + \text{Non-lin\'earit\'e} +
    \text{Pooling}}
\end{equation}

\noindent Apr\`es plusieurs blocs, les cartes de features sont
aplaties et connect\'ees \`a des couches enti\`erement connect\'ees
(fully connected) pour la classification finale.

\begin{important}
  Les ConvNets exploitent la structure locale des donn\'ees~:
  pixel voisins, \'echantillons temporels proches, etc. Chaque
  bloc extrait des motifs de plus en plus abstraits.
\end{important}

% ===================================================================
\section{\'Etude de cas~: LeNet-5 (LeCun et al., 1989)}
% ===================================================================

\subsection{Architecture}

LeNet-5 est le premier r\'eseau convolutionnel moderne, d\'evelopp\'e
pour la reconnaissance de chiffres manuscrits (MNIST).

\begin{center}
\begin{tabular}{clccc}
\toprule
Couche & Type & Taille & Maps $\times$ Dim. & \# Param. \\
\midrule
Entr\'ee & Image & $32\times 32$ & $1 \times 32 \times 32$ & 0 \\
C1 & Convolution ($5\times5$, 6 filtres) & & $6 \times 28 \times 28$ & 156 \\
S2 & Pooling ($2\times2$, stride 2) & & $6 \times 14 \times 14$ & 12 \\
C3 & Convolution ($5\times5$, 16 filtres) & & $16 \times 10 \times 10$ & 1516 \\
S4 & Pooling ($2\times2$, stride 2) & & $16 \times 5 \times 5$ & 32 \\
C5 & Convolution ($5\times5$, 120 filtres) & & $120 \times 1 \times 1$ & 48\,210 \\
F6 & Fully connected & & 84 & 10\,164 \\
F7 & Fully connected (sortie) & & 10 & 850 \\
\midrule
\multicolumn{4}{r}{Total} & $\sim$60\,000 \\
\bottomrule
\end{tabular}
\end{center}

\subsection{D\'etails}

\begin{description}
  \item[C1]~: $6$ filtres $5\times5$, pas de padding $\Rightarrow$
    $28\times28$. Param\`etres~: $(5\times5 + 1)\times 6 = 156$.
  \item[S2]~: sous-\'echantillonnage $2\times2$ stride $2$, somme
    pond\'er\'ee + biais $\Rightarrow$ $14\times14$.
  \item[C3]~: $16$ filtres $5\times5$, connexion s\'lective aux
    cartes de S2 (non fully connected) $\Rightarrow$ $10\times10$.
  \item[S4]~: identique \`a S2 $\Rightarrow$ $5\times5$.
  \item[C5]~: $120$ filtres $5\times5\times16$ (profondeur totale)
    $\Rightarrow$ cartes $1\times1$ (vecteur de $120$).
  \item[F6]~: $84$ neurones fully connected.
  \item[F7]~: $10$ neurones de sortie (softmax).
\end{description}

\subsection{R\'esultats sur MNIST}

\begin{itemize}
  \item Ensemble original~: $60\,000$ images, erreur test~$0{,}95\%$.
  \item Avec $540\,000$ distortions artificielles~: erreur test
    $0{,}8\%$.
  \item D\'eploiement r\'eussi pour la lecture de codes postaux
    aux \'Etats-Unis.
\end{itemize}

% ===================================================================
\newpage
\section*{Questions cl\'es (QCM)}
% ===================================================================
\addcontentsline{toc}{section}{Questions cl\'es (QCM)}

\begin{enumerate}

  \item Quelle est la principale limitation des MLP pour les
    images~?
    \begin{enumerate}
      \item Ils sont trop lents \`a l'inf\'erence
      \item Explosion du nombre de param\`etres et ignorance de la
        topologie spatiale
      \item Ils ne peuvent pas apprendre
      \item Ils n\'ecessitent trop de donn\'ees
    \end{enumerate}
    \textbf{R\'eponse~: b (explosion des param\`etres et perte de
    la topologie).}

  \item Que permet le partage de poids (shared weights)~?
    \begin{enumerate}
      \item Augmenter le nombre de param\`etres
      \item R\'eduire consid\'erablement le nombre de param\`etres
        \`a apprendre
      \item Acc\'el\'erer le calcul
      \item \'Eviter le sur-apprentissage
    \end{enumerate}
    \textbf{R\'eponse~: b (r\'eduction du nombre de param\`etres).}

  \item Quelle est la diff\'erence entre convolution et
    corr\'elation crois\'ee~?
    \begin{enumerate}
      \item Aucune
      \item La corr\'elation crois\'ee ne sym\'etrise pas le masque
      \item La convolution est plus rapide
      \item La corr\'elation crois\'ee utilise des filtres plus
        grands
    \end{enumerate}
    \textbf{R\'eponse~: b (pas de sym\'etrisation du masque).}

  \item Quel est le r\^ole d'une couche de pooling~?
    \begin{enumerate}
      \item Augmenter la r\'esolution
      \item Apporter une invariance aux translations locales et
        r\'eduire la dimension spatiale
      \item Remplacer la convolution
      \item Initialiser les poids
    \end{enumerate}
    \textbf{R\'eponse~: b (invariance aux translations et
    sous-\'echantillonnage).}

  \item Dans LeNet-5, combien de filtres la couche C1
    poss\`ede-t-elle~?
    \begin{enumerate}
      \item 1
      \item 6
      \item 16
      \item 120
    \end{enumerate}
    \textbf{R\'eponse~: b (6 filtres $5\times5$).}

  \item Quelle est la taille des cartes de sortie de C1 dans
    LeNet-5~?
    \begin{enumerate}
      \item $32\times32$
      \item $28\times28$
      \item $14\times14$
      \item $10\times10$
    \end{enumerate}
    \textbf{R\'eponse~: b ($28\times28$, convolution $5\times5$
    sans padding).}

  \item Combien de param\`etres LeNet-5 a-t-il au total~?
    \begin{enumerate}
      \item $\sim$6\,000
      \item $\sim$60\,000
      \item $\sim$600\,000
      \item $\sim$6\,000\,000
    \end{enumerate}
    \textbf{R\'eponse~: b ($\sim$60\,000 param\`etres).}

  \item Quel jeu de donn\'ees a servi pour \'evaluer LeNet-5~?
    \begin{enumerate}
      \item ImageNet
      \item MNIST
      \item CIFAR-10
      \item COCO
    \end{enumerate}
    \textbf{R\'eponse~: b (MNIST, chiffres manuscrits).}

  \item Quelle performance LeNet-5 atteint-elle sur MNIST avec
    distortions artificielles~?
    \begin{enumerate}
      \item $5\%$ d'erreur
      \item $0{,}8\%$ d'erreur
      \item $10\%$ d'erreur
      \item $0{,}1\%$ d'erreur
    \end{enumerate}
    \textbf{R\'eponse~: b ($0{,}8\%$ d'erreur test).}

  \item Qui a propos\'e LeNet-5~?
    \begin{enumerate}
      \item Hinton et al.
      \item LeCun et al.
      \item Krizhevsky et al.
      \item Goodfellow et al.
    \end{enumerate}
    \textbf{R\'eponse~: b (LeCun et al., 1989).}

\end{enumerate}

% ===================================================================
\begin{thebibliography}{9}
% ===================================================================

\bibitem{LeCun1989}
  Y.~LeCun, B.~Boser, J.~S.~Denker, D.~Henderson, R.~E.~Howard,
  W.~Hubbard, L.~D.~Jackel.
  \newblock \emph{Backpropagation applied to handwritten zip code
    recognition}.
  \newblock Neural Computation, 1(4):541--551, 1989.

\bibitem{Goodfellow2016}
  I.~Goodfellow, Y.~Bengio, A.~Courville.
  \newblock \emph{Deep Learning}.
  \newblock MIT Press, 2016.

\end{thebibliography}

\end{document}
